\section{Experimental Setup}

\subsection{Implementation}

Experiments use Python~\R{PythonVersion} and PyTorch~\R{TorchVersion} on
\R{GPUName}. All reported environment details are emitted automatically by the
benchmarking script rather than transcribed by hand. Code, configurations,
checkpoints and the complete run log are released.

\subsection{Datasets}

\textbf{PTB-XL}~\cite{wagner2020,strodthoff2021} (v1.0.3) provides 21{,}799
twelve-lead resting ECGs of 10\,s duration from 18{,}885 patients, released
under CC-BY-4.0. We use the 100\,Hz variant and the five superdiagnostic
classes: NORM (44.5\%), MI (25.6\%), STTC (24.5\%), CD (22.9\%) and HYP
(12.4\%); approximately 28\% of recordings carry multiple labels. Following
the official protocol, folds 1--8 form the training set, fold 9 validation and
fold 10 the held-out test set. After removing recordings without diagnostic
labels, 21{,}388 remain.

\textbf{External corpora.} For cross-dataset evaluation we use CPSC-2018,
Chapman--Shaoxing/Ningbo, and the Georgia twelve-lead database, all from the
PhysioNet/CinC-2020 collection. Source annotations use SNOMED-CT codes with a
finer, more rhythm-oriented vocabulary than PTB-XL's diagnostic superclasses;
our mapping is many-to-one and necessarily lossy, and records whose codes map
to no superclass are dropped. These results therefore measure transfer of the
\emph{superclass concept} under distribution shift, not identical-task
performance, and should not be read as directly comparable to the PTB-XL
numbers.

\subsection{Baselines}

\textbf{(1) Fixed-dimension, same backbone.} For each backbone and each $m \in
\mathcal{M}$, a separate encoder with a single classifier at dimension $m$,
trained with identical data, augmentation, schedule, and seeds. This is the
comparison that isolates the effect of nesting. The earlier version of this
work compared an Inception1D MRL model against XResNet1D-101 fixed baselines
and reported the difference as evidence for MRL; that comparison confounds
nesting with backbone choice and cannot support the conclusion. Both
confounds are separated here.

\textbf{(2) SVD compression.} Truncated SVD of the 512-dimensional embeddings
with logistic-regression classifiers on the projected features. Both the
basis and the classifier are fitted on the \emph{training} split and applied
unchanged to the test split. The earlier version fitted both on the test
partition and evaluated on the same partition; although labelled
``optimistically biased'', it was nonetheless plotted as an upper bound
against honestly evaluated models.

\textbf{(3) MRL-E.} The shared-weight variant of Equation~\ref{eq:mrle}.

\textbf{(4) Published PTB-XL systems.} Reported for context in
Section~\ref{sec:sota}, with the explicit caveat that these are
non-adaptive systems evaluated on a different axis.

\subsection{Hyperparameters}

AdamW, learning rate $10^{-3}$, weight decay $10^{-2}$, batch size 128, cosine
annealing with 5-epoch linear warmup, up to 80 epochs with early stopping
(patience 15) on mean validation macro AUC. Label smoothing is disabled by
default and ablated separately. All hyperparameters are fixed on fold 9; no
tuning is performed on fold 10.

\subsection{Statistical Protocol}
\label{sec:statproto}

Every model is trained with five seeds ($\{0,\dots,4\}$); tables report mean
$\pm$ standard deviation. Test-set uncertainty is quantified by bootstrap over
recordings ($B$=2{,}000).

For comparisons between models evaluated on the same test set we use a
\emph{paired} bootstrap resampling identical recording indices for both
models, and DeLong's test~\cite{delong1988,sun2014} per class. Comparisons
across the six nesting dimensions are corrected by the Holm step-down
procedure.

The central claim---that performance is stable across dimensions---is a claim
of \emph{no meaningful difference}, which a superiority test cannot establish;
failing to reject a null hypothesis is not evidence for it. We therefore
pre-register an equivalence margin of $\pm$0.01 macro AUC and require the
entire 95\% confidence interval of the difference to fall inside that margin.
The margin is chosen as roughly half the typical seed-to-seed standard
deviation of a PTB-XL classifier, so that ``equivalent'' means a difference
smaller than the noise floor of the training procedure itself.

\subsection{Evaluation Metrics}

The primary metric is macro-averaged AUC-ROC following the PTB-XL
standard~\cite{strodthoff2021}, which weights the rare HYP class equally.
Secondary metrics are macro F1 (per-class thresholds tuned on validation) and
macro average precision. All test metrics are computed on fold 10
($n$=2{,}158).
